\chapter{Implementazione}
\label{chap:implementazione}

% ==================================================
% 5.1 Organizzazione del codice
% ==================================================

\section{Organizzazione del codice applicativo}
\label{sec:organizzazione-codice}

% Obiettivo:
% mostrare come l'architettura del Capitolo 4 viene
% realizzata nel repository.
%
% Non trasformare la sezione in una lista di file: spiegare
% le responsabilità implementative e le dipendenze tra moduli.

% ==================================================
% 5.2 Autenticazione e profilo
% ==================================================

\section{Autenticazione, conferma email e profilo utente}
\label{sec:implementazione-auth-profilo}

% Descrivere il flusso reale:
% - registrazione;
% - hash della password;
% - token di conferma email;
% - login con sessione;
% - protezione delle route;
% - gestione profilo e preferenza di saluto.

% ==================================================
% 5.3 Catalogazione libraria
% ==================================================

\section{Catalogazione libraria e gestione delle immagini}
\label{sec:implementazione-catalogazione}

% Contenuti attesi:
% - creazione, modifica ed eliminazione libri;
% - persistenza tramite repository e Prisma;
% - gestione di copertine, immagini e thumbnail;
% - visibilità e disponibilità;
% - statistiche di visualizzazione.

% ==================================================
% 5.4 Catalogazione assistita
% ==================================================

\section{Catalogazione assistita tramite ISBN, OCR e metadati esterni}
\label{sec:implementazione-catalogazione-assistita}

% Derivare la descrizione dal codice in packages/ai e dalle
% server action della feature libri.
%
% Principio da evidenziare:
% il sistema propone dati, ma il salvataggio resta sotto
% controllo esplicito dell'utente.

% ==================================================
% 5.5 Geolocalizzazione e ricerca nearby
% ==================================================

\section{Geolocalizzazione, query spaziali e visualizzazione cartografica}
\label{sec:implementazione-geolocalizzazione}

% Descrivere:
% - geocoding e autocomplete;
% - salvataggio di coordinate private/pubbliche;
% - query PostGIS per raggio e distanza;
% - fallback in TypeScript se applicabile;
% - MapLibre, marker, cluster e adattamenti mobile.

% ==================================================
% 5.6 Richieste e dashboard
% ==================================================

\section{Richieste tra utenti, dashboard e amministrazione}
\label{sec:implementazione-richieste-dashboard}

% Collegare le funzionalità implementate ai requisiti:
% - invio richiesta;
% - vincoli su proprietario/autenticazione/disponibilità;
% - gestione delle richieste ricevute e inviate;
% - dashboard utente;
% - dashboard amministrativa.
